% Generated draft: requires mathematical and source review.
\documentclass[11pt]{article}
\usepackage[margin=25mm]{geometry}
\usepackage{fontspec}
\setmainfont{Latin Modern Roman}
\usepackage{amsmath,amssymb,mathtools,unicode-math}
\setmathfont{Latin Modern Math}
\usepackage{enumitem,needspace,xurl,hyperref}
\setlength{\parindent}{0pt}
\setlength{\parskip}{5pt}
\setlength{\emergencystretch}{3em}
\newcommand{\N}{\mathbb{N}}
\newcommand{\Z}{\mathbb{Z}}
\newcommand{\Q}{\mathbb{Q}}
\newcommand{\R}{\mathbb{R}}
\newcommand{\C}{\mathbb{C}}
\newcommand{\F}{\mathbb{F}}
\newcommand{\A}{\mathbb{A}}
\newcommand{\PP}{\mathbb{P}}
\newcommand{\E}{\mathbb{E}}
\newcommand{\eps}{\varepsilon}
\newcommand{\dd}{\,\mathrm{d}}
\begin{document}
\section*{UnsolvedMath ID: 30005969}
\textbf{Source problem code:} catalogue entry
\subsection{Definitions and mathematical statement}
The supplied record asks an empirical question about individual extreme rainfall events in the northeastern United States. It does not specify a mathematical conjecture or a complete statistical model.

For schematic notation only, let \(\lambda(t)\) denote the frequency of individual extreme rainfall events at observation time or period \(t\), and let \(F_t\) denote the distribution of their individual intensities. These symbols require a specified observation region, observation periods, event identification rule, intensity measure, and sampling or probability model; none is fully supplied. In particular, no Poisson assumption, independence assumption, or particular ordering of the distributions \(F_t\) is asserted.

\textbf{Empirical question.} Is climate change making individual extreme rainfall events in the northeastern United States more frequent, more intense, or both? In the schematic notation, the requested investigation concerns increases in \(\lambda(t)\), changes in \(F_t\) corresponding to greater event intensity, and whether those changes can be attributed to climate change. The source does not specify the comparison periods, the meaning of an intensity increase, or an attribution criterion. Consequently, this question cannot be converted into a unique quantified mathematical assertion from the supplied evidence.

\textbf{Supplied qualifications.} The registry labels the statement \texttt{unrecoverable} and its status \texttt{partially\_solved}. Its supplied statement-verification note says that the official report poses this empirical question and then reports that events became more frequent, while their individual-size distribution showed little evidence of increase. This is a supplied account of an empirical finding, not an independently verified theorem; little evidence of an increase does not establish that the intensity distribution is unchanged.
\subsection{Short English statement}
Determine whether climate change increases the frequency, individual intensity, or both, of extreme rainfall events in the northeastern United States. The supplied account reports increased frequency and little evidence of an intensity-distribution increase, with broader attribution remaining scope-dependent. Event definitions, observation periods, statistical targets, and attribution assumptions are missing, so a unique precise mathematical conjecture cannot be recovered.
\subsection{Sources}
\begin{itemize}
\item Local supplied UnsolvedMath registry record, permanent ID \(30005969\), title \emph{Observed Extreme-Rainfall Trends}, source problem code \texttt{catalogue entry}, collection \emph{Oberwolfach Reports Open Problems}. The supplied proposed year is 2024; the registry labels the statement \texttt{unrecoverable} and the status \texttt{partially\_solved}.
\item Supplied report citation: \emph{Mathematics, Statistics, and Geometry of Extreme Events in High Dimensions} (2025), Oberwolfach Reports, DOI 10.4171/owr/2024/38, pp.\ 2179--2250. Original OWR report: \url{https://doi.org/10.4171/owr/2024/38}. The background describes these reports as documenting mathematical workshops and their open-problem discussions.
\item Supplied presentation attribution: \emph{Rainfall trends in the USA: data quality and prediction skill}, Holger Rootz\'{e}n, joint work with Helga Olafsdottir, David Bolin, and Richard Katz. The registry does not separately identify a proposer.
\item Supplied literature source, labeled ``Olafsdottir--Rootzen--Bolin primary study'': \url{https://journals.ametsoc.org/view/journals/clim/34/22/JCLI-D-20-0938.1.xml}.
\item Literature assessment dated exactly \texttt{2026-08-22}: the supplied curation describes a peer-reviewed study of 333 northeastern U.S. gauges as finding temperature-linked increases in event frequency and little evidence of an intensity-distribution trend; broader attribution remains scope-dependent. This assessment and the partially solved label are supplied dated curation derived from the OWR material, not independently verified here. The supplied background explicitly requests independent verification of time-sensitive status and cited claims.
\end{itemize}
\end{document}
